\section{Experiments}

\subsection{Experimental Setup}
\paragraph{Datasets.}
We train on the referring video object segmentation training split and evaluate on validation splits.
For each expression, the model ranks all recalled candidates from the same video-expression pair.

\paragraph{Implementation Details.}
We use Qwen3-VL as the base VLM and apply LoRA fine-tuning.
The visual encoder is frozen, while LoRA adapters are injected into language-side attention and MLP projections.
Training is performed with listwise candidate groups (batch size $=1$ group), AdamW optimizer, and mixed precision.
The candidate masks used for supervision are generated in an offline preprocessing stage.

\paragraph{Input Construction.}
Each candidate is represented by a sampled multi-frame tube.
For each frame, we overlay candidate mask (or bbox) on the original image and pair it with a unified Yes/No prompt.
The model outputs $P(\text{Yes})$ for each candidate, which is used for ranking.

\paragraph{Supervision.}
Candidate-level IoU against GT is precomputed offline.
Positives/negatives are defined by IoU thresholds $\tau_{pos}$ and $\tau_{neg}$.
We optimize listwise ranking loss and soft mask-aware alignment loss jointly.

\paragraph{Evaluation Metrics.}
We report region similarity $\mathcal{J}$ (IoU), contour accuracy $\mathcal{F}$, and their mean $\mathcal{J}\&\mathcal{F}$.
Unless otherwise stated, expression-level averages are used as the primary comparison metric.

\subsection{Main Results}
\begin{table}[t]
\centering
\caption{Main comparison on the validation split.}
\begin{tabular}{lccc}
\hline
Method & $\mathcal{J}$ & $\mathcal{F}$ & $\mathcal{J}\&\mathcal{F}$ \\
\hline
Pretrained Qwen3-VL + CRF-SAM3 & -- & -- & -- \\
Ours (Listwise + Soft Alignment + LoRA) & -- & -- & -- \\
\hline
\end{tabular}
\label{tab:main_results}
\end{table}

\subsection{Ablation Study}
\begin{table}[t]
\centering
\caption{Ablation of key components.}
\begin{tabular}{lccc}
\hline
Configuration & $\mathcal{J}$ & $\mathcal{F}$ & $\mathcal{J}\&\mathcal{F}$ \\
\hline
Pointwise supervision only & -- & -- & -- \\
+ Listwise ranking & -- & -- & -- \\
+ Soft mask-aware alignment & -- & -- & -- \\
+ LoRA adaptation (full setting) & -- & -- & -- \\
\hline
\end{tabular}
\label{tab:ablation}
\end{table}

\subsection{Analysis}
\paragraph{Why Listwise Helps.}
Listwise optimization aligns training objective with inference-time ranking among all candidates under the same expression, reducing score calibration mismatch from independent pointwise training.

\paragraph{Why Soft Alignment Helps.}
The soft alignment loss increases attention/feature probability mass inside target masks, suppressing background shortcut cues and improving instance-level discrimination.

\paragraph{Efficiency.}
Offline cache construction decouples expensive candidate recall from VLM fine-tuning, enabling stable and memory-efficient training iterations.

\subsection{Reproducibility Notes}
\begin{itemize}
\item Keep train/eval prompt template and scoring function strictly aligned (Yes/No probability scoring).
\item Ensure full-frame prediction coverage for each expression during evaluation.
\item Use identical candidate-generation and ranking configuration when comparing pretrained vs fine-tuned models.
\end{itemize}
